\section{引言}

给定实体对及其图上下文，关系预测需要从关系词表中识别相应的关系。GRIP 将图结构融入语言模型参数，通过生成式问答完成这一任务 \citep{feng2025grip}。在原有生成目标之外加入候选级对比损失，可以让目标关系与替代关系直接竞争；这一训练信号的作用取决于替代关系如何选取。

候选过于容易时，提供的排序边界信号可能有限；只按模型高分挑选时，候选又可能反映偶然的分数混淆。更关键的是，知识图谱通常不完整：$(h,r,t)$ 未被标注，并不等于该事实为假。困难候选因此同时面临\emph{信息量}和\emph{事实不确定性}两项要求。本文聚焦一个可检验的问题：能否在保持 GRIP 原始生成路径的同时，构造来源可追溯、排除训练已知关系的候选库，并用它训练关系排序边界？

为此，我们提出结构反事实对比学习（\scc）。首先，在训练关系词表内排除正关系及同一实体对的 train-known 关系；随后，用训练图的关系画像衡量结构相似度，用冻结的 generation-only 基线 \bone{} 衡量候选分数违例度。两种信号经带均匀下限的分布组合，每题无放回采样九个候选，并在 Stage~2 前冻结。训练仅为原生成目标增加候选级 InfoNCE，不改变 Stage~1 或原始提示。这里的“结构反事实”指对已观测实体对替换关系标签，不表示剩余候选都是真实负例。图~\ref{fig:mechanism} 展示训练流程。

NELL23K 的一次 seed=2026 运行完成了候选库、Stage~2 训练和全测试评估，在 4,943 条问题上得到官方 EM 91.52\% 和 MRR 95.27\%。GRIP 原论文报告的 NELL23K EM 为 87.74\%；按论文主指标计算，\scc{} 高出 3.78 个百分点。在相同候选列表的前 200 条问题上，\scc{} 与本地 \bone{} 的 MRR 分别为 94.96\% 和 93.36\%。FB15K237-2 的历史候选对比训练实例也取得官方 EM 91.62\%，高于原论文的 88.75\%；这一路径使用不同候选文件和训练配置，作为辅助观察，而非 SC-CL 的跨图复现。

本文的具体贡献是：定义同时利用结构相似度和冻结模型违例度的训练期候选分布；将 train-known 过滤、固定候选库和来源记录落实为可审计约束；在官方关系预测设置下报告 EM、MRR 与 Hits，并给出历史候选训练实例的辅助结果。全文限于固定实体对的关系预测。
